% Drawing 74/02-5 -- folder 74, batch 02, archive page 26 (PDF page 27).
% Transcription: une petite figure en angle, puis deux étoiles de segments reliées par une arête ; au-dessous un faisceau de traits, une « maison » (carré surmonté d'un toit, un triangle inscrit, un axe pointillé), un segment divisé en trois, un graphe biparti de trois sommets en haut et trois en bas reliés par de nombreuses arêtes, et plusieurs petits schémas de trois traits horizontaux reliés par des obliques, l'un d'eux étiqueté $\alpha$, $\beta$, $\alpha$.
% Kind: page of small graph sketches. Method 2 (schematic redraw from the facsimile). Status: partial.
% Notes: Every stroke redrawn where it lies on the page. A small angle (a vertical stroke and an oblique from its top). Two stars, each made of thin wedges (two strokes from the centre closed by a short cap or an arc), four wedges in the left star and five in the right one, their centres joined by a thick curve; a bold point below them, joined to the left centre by a curve and to the right one by a stroke. Below the point, a fan of five strokes that start about two units under it, not at it. The house: the left roof stroke rising from the top of the left wall, the right roof stroke starting just beside its top and ending in a hook, the eave line, the left and right walls (the right one standing apart from the roof), the base, the triangle inscribed from the foot corners to a point on the axis, and the dashed axis from under the apex down to that point. A segment of three pieces between brackets. The graph of three columns: a top bar (doubled at the left), the three columns, a V from the top left down to the middle column and up to the top right, and the crossing strokes between the columns, with blots at the column points. The small schemas of horizontal strokes joined by obliques (one with alpha, beta, alpha under it, alpha with a dot over it), the strokes with an arrow at the lower middle, and a Z-shaped schema at the lower left; their strokes do not meet at points, as on the page. The arithmetic and text around (9 x 9 = 81, 3x3x3 = 27, (40), 54, ...) are text and not drawn. Labels at \small, the figure drawn at 1.5 mm per unit so that they keep their places. Redrawn 2026-11-09 (issue 128); overlay check 2026-11-13 (issue 147). Second pass 2026-11-14 (issue 147, Claude claude-opus-5-5): every stroke re-read on 14 to 40 px/unit gridded crops of the 400 dpi colour scan and redrawn there (the stars were drawn as single spokes and the fan from the point; the small schemas lay 2 to 3 units off). Rasterised by the session's previewer (not compiled: no TeX in the session) and laid over the crop: every stroke lies on its ink to within about 1 unit. Still differs: in the graph of three columns the crossing strokes are dense, doubled and overinked; the eight drawn (three from the left column to the top of the middle one, one each from the top of the middle column to the lower right and to the right column, one from the top right to the middle column, and the V) are a reading of their number and of where each ends, and the transcription's « graphe biparti ... reliés par de nombreuses arêtes » is not checked as a complete bipartite graph. Not checked by a reader.
% Coordinates: 1 unit = 1.5 mm = 1/100 of the crop width (bbox in page fractions: [0.42, 0.43, 0.9, 0.86]).
% Redrawn by Claude (claude-opus-5-5), issue #73, 2026-11-09; not checked by a human; not compiled (no TeX in the session).
% Checked by Claude (claude-opus-5-5) 2026-11-13 (issue #147): overlaid on the 300 dpi colour facsimile; not compiled (no TeX in the session). Status: partial.
% Second pass by Claude (claude-opus-5-5) 2026-11-14 (issue #147): redrawn stroke by stroke from the 400 dpi colour facsimile, overlaid on a grid; not compiled (no TeX in the session). Status: partial.
% Rights: reproduces part of the work -- behind RENDER_TRANSCRIPTS, never published by itself.
\documentclass[tikz,border=3pt]{standalone}
\usepackage{amsmath,amssymb}
\usetikzlibrary{patterns,arrows.meta}
\tikzset{pencil/.style={black!45}, struck/.style={black!35}, mark/.style={circle,draw,inner sep=1.3pt}, pt/.style={circle,fill,inner sep=0.9pt}, opt/.style={circle,draw,fill=white,inner sep=1.1pt}}
\begin{document}
\begin{tikzpicture}[x=1.5mm,y=1.5mm,line cap=round,line join=round,font=\scriptsize]
\draw[thick] (20.5,-21.5) -- (20.2,-11.3) -- (27.1,-20);
% left star: four wedges from its centre
\draw (36.3,-20.8) -- (36.2,-14.9) -- (38.9,-15.1);
\draw (38.4,-15.3) -- (36.3,-20.8);
\draw plot[smooth] coordinates {(36.3,-20.8) (35.1,-18.1) (33.8,-16.1) (32.8,-15.2)};
\draw plot[smooth] coordinates {(32.8,-15.2) (31.6,-16.6) (30.5,-18.4)};
\draw (30.2,-18.6) -- (36.3,-20.8);
\draw (30.2,-18.6) -- (30.8,-19.3);
\draw (36.3,-20.8) -- (30.5,-22.5);
\draw plot[smooth] coordinates {(30.9,-23.1) (30.8,-25.4) (31.1,-27.0)};
\draw (31.1,-27.0) -- (36.3,-20.8);
\draw (36.3,-20.8) -- (34,-27.4) -- (36.2,-27.2);
\draw (36.6,-26.6) -- (36.3,-20.8);
\fill (37.4,-27.2) circle[radius=0.15];
% right star
\draw (43.2,-13.1) -- (45.3,-18.2);
\draw (44.4,-12.9) -- (45.9,-12.3) -- (46,-18.4);
\draw (46,-18.4) -- (48.5,-12.3);
\draw (49.4,-12.3) -- (50.6,-14.8);
\draw (50.6,-15.1) -- (46,-18.4);
\draw (46,-18.4) -- (53.2,-16.2) -- (54,-16.3) -- (50.7,-20.5) -- (46.4,-18.8);
\draw (46,-18.4) -- (47.9,-22.9) -- (49.9,-21.7) -- (46.4,-18.8);
% the two centres joined; the bold point below them
\draw[thick] plot[smooth] coordinates {(36.3,-20.8) (38.9,-20.4) (42.4,-19.3) (46,-18.4)};
\draw plot[smooth] coordinates {(36.8,-21.2) (39.4,-22.6) (41.7,-24.5) (42.2,-25.4)};
\draw (42.4,-25.2) -- (45.6,-19.6);
\fill (42.3,-25.6) circle[radius=0.45];
% the fan, its strokes starting about two units below the point
\draw (40.7,-27.5) -- (32.7,-35.8);
\draw (41.2,-27.3) -- (37.7,-34.8);
\draw plot[smooth] coordinates {(42.2,-28) (42.6,-30.7) (43.5,-34.7)};
\draw (43,-27.7) -- (47.8,-34.5);
\draw plot[smooth] coordinates {(44.3,-27.5) (47.7,-28.7) (49.7,-30.7) (51.2,-32)};
% the house
\draw (31.3,-43) -- (35.7,-38.7) -- (41,-34);
\draw plot[smooth] coordinates {(41.3,-34.7) (43,-35) (45.7,-37.3) (48,-39.7)};
\draw (48,-39.7) -- (48.2,-39.4);
\draw plot[smooth] coordinates {(31.3,-43) (35.7,-41.7) (40.7,-41) (46,-40.3)};
\draw (46,-40.3) -- (45.3,-40.7);
\draw (31.3,-43) -- (31.2,-53.7) -- (31.6,-53.9);
\draw plot[smooth] coordinates {(47,-41.8) (47.2,-46) (47.2,-50) (47.8,-51.3)};
\draw (31.5,-55.2) -- (46,-53.5) -- (40.8,-45.7) -- cycle;
\draw[densely dashed] (40.7,-35.5) -- (40.7,-45.2);
% the segment: three pieces between brackets
\draw[thick] (39.1,-65.1) -- (47.6,-64.8);
\draw[thick] (50.1,-64.9) -- (59.5,-64.2);
\draw[very thick] (60.9,-64.2) -- (70.6,-63.8);
\draw (39.1,-64.2) -- (39.1,-66.9);
\draw (47.6,-64) -- (47.6,-65.6);
\draw (50.1,-64) -- (50.1,-65.6);
\draw (59.5,-63.3) -- (59.5,-65);
\draw (60.9,-63.2) -- (60.9,-65);
\draw[thick] (70.6,-62.8) -- (70.6,-64.9);
% the schema of three columns
\draw (32.6,-77) -- (41,-76);
\draw[thick] (32.6,-77.4) -- (41.8,-76.5) -- (50.1,-77.1);
\draw (42.3,-75.8) -- (42.8,-76.5);
\draw (33.8,-77.8) -- (33.9,-87.8);
\draw (41.8,-76.8) -- (41.8,-87.8);
\draw (49.5,-77.5) -- (49.6,-87.8);
\draw plot[smooth] coordinates {(33.8,-78) (37.2,-80.5) (41.6,-82.4)};
\draw (41.6,-82.4) -- (49,-77.8);
\draw (33.5,-80.8) -- (41.5,-76.5);
\draw (33.5,-84.8) -- (40.2,-76.6);
\draw (34,-85.6) -- (41,-76.8);
\draw (42,-76.8) -- (49.5,-86.8);
\draw (42.5,-76.5) -- (48.5,-81);
\draw (49,-77.8) -- (42,-85.6);
\fill (33.4,-77.5) circle[radius=0.4];
\fill (33.5,-82.5) circle[radius=0.4];
\fill (33.4,-85.8) circle[radius=0.4];
\fill (41.8,-85.8) circle[radius=0.4];
\fill (49,-81.4) circle[radius=0.4];
\fill (49.3,-86.8) circle[radius=0.4];
\fill (49.6,-77) circle[radius=0.4];
% the Z-like schemas at the right
\draw (61.4,-76.5) -- (62.3,-76.5);
\draw (63.4,-76.8) -- (69.4,-76.2);
\draw (70.9,-76) -- (71.4,-76.5);
\draw (73.1,-76.3) -- (81.6,-76);
\draw[very thick] (74.4,-76.7) -- (81.6,-76.1);
\draw[thick] (63.8,-83.2) -- (70.3,-77.4);
\draw (66.3,-80.3) -- (68.9,-77.3) -- (68.4,-78.2);
\draw (62.3,-83.3) -- (63.1,-83.2);
\draw (64.7,-83.1) -- (69.4,-82.3);
\draw[thick] (64.4,-87.8) -- (72.2,-81.8) -- (72.1,-81.5);
\draw (64.2,-88.5) -- (70,-84);
\draw (62.3,-88.5) -- (62.8,-89.3);
\draw (71.9,-87.4) -- (72.2,-88.4);
\draw[very thick] (73.4,-88.5) -- (81.1,-87.7);
\draw (81.3,-86.5) -- (81.6,-87.4);
\draw (73.3,-77.6) -- (73.9,-78.1) -- (80.6,-81.1);
\draw[thick] (74.7,-79) -- (78.1,-79.6);
\draw (81.6,-81) -- (82,-81);
\draw (74.4,-81.8) -- (80,-81.8);
% lower schemas
\draw (49.9,-100.5) -- (54.8,-102.7);
\draw (52.5,-105.6) -- (55.7,-104);
\draw[->] (52.7,-107.9) -- (59.5,-107);
\draw (54.3,-113.4) -- (59.8,-109.7);
\draw[very thick] (73.2,-100.5) -- (80.5,-100);
\draw[thick] (80.5,-100.1) -- (73.9,-107);
\draw (75.5,-104) -- (79.5,-100.9);
\draw (81.8,-100.2) -- (87.5,-100.2);
\draw[pencil] (77.3,-104.7) -- (79.1,-107.5);
\draw (79.8,-106.1) -- (87.7,-105.9);
\draw (80.7,-112.2) -- (87,-111.9);
\fill (74.8,-113.1) circle[radius=0.2];
\node[inner sep=1pt] at (74.8,-116.1) {\small $\alpha$};
\node[inner sep=1pt] at (81.4,-116.5) {\small $\beta$};
\node[inner sep=1pt] at (87,-116.3) {\small $\alpha$};
\draw (27.75,-114.9) -- (33,-114.7);
\draw (32.3,-114.8) -- (27.6,-118.4);
\draw[very thick] (31.9,-117.8) -- (35.9,-117.7);
\fill (36.3,-115) circle[radius=0.2];
\draw[very thick] (27.75,-121) -- (32.4,-120.4);
\draw (35.75,-120.4) -- (35.9,-121);
\end{tikzpicture}
\end{document}
